% SEGMENT OLP-0014-S01
% Part: sets-functions-relations
% Chapter: relations
% Section: special-properties

\documentclass[../../../include/open-logic-section]{subfiles}

\begin{document}

\olfileid{sfr}{rel}{prp}
\olsection{தொடர்புகளின் சிறப்புப் பண்புகள்}

\begin{intro}
சில வகைத் தொடர்புகள் மிகவும் பொதுவாகப் பயன்படுவதால்,
அவற்றுக்குச் சிறப்புப் பெயர்கள் வழங்கப்பட்டுள்ளன.
எடுத்துக்காட்டாக, $\le$, $\subseteq$ ஆகிய இரண்டும் தத்தமது
சார்பகங்களை ஒத்த முறைகளில் தொடர்புபடுத்துகின்றன
($\le$ க்கு $\Nat$ என்பதையும், $\subseteq$ க்கு $\Pow{A}$ என்பதையும்
சார்பகங்களாகக் கொள்ளலாம்).
இத்தொடர்புகள் எவ்வாறு ஒத்திருக்கின்றன, எவ்வாறு வேறுபடுகின்றன என்பதைத்
துல்லியமாக அறிய, தொடர்புகள் கொண்டிருக்கக்கூடிய சில சிறப்புப் பண்புகளின்படி
அவற்றை வகைப்படுத்துகிறோம். இந்தச் சிறப்புப் பண்புகளில் சிலவும் அவற்றின்
சேர்க்கைகளும் குறிப்பாக முக்கியமானவையாகின்றன:
வரிசைகளும் சமானத் தொடர்புகளும் அவற்றில் அடங்கும்.
\end{intro}

% SEGMENT OLP-0014-S02
\begin{defn}[தற்சுட்டுப் பண்பு]
ஒவ்வொரு $x \in A$ க்கும் $Rxx$ ஆக இருக்கும்போது,
அப்போதும் அப்போது மட்டுமே $R \subseteq A^2$ என்னும் தொடர்பு
\emph{தற்சுட்டுப் பண்புடையது} எனப்படும்.
\end{defn}

% SEGMENT OLP-0014-S03
\begin{defn}[கடப்புப் பண்பு]
$Rxy$, $Ryz$ ஆகியவை உண்மையாக இருக்கும்போதெல்லாம் $Rxz$ என்பதும்
உண்மையாக இருக்கும்போது, அப்போதும் அப்போது மட்டுமே
$R \subseteq A^2$ என்னும் தொடர்பு \emph{கடப்புப் பண்புடையது} எனப்படும்.
\end{defn}

% SEGMENT OLP-0014-S04
\begin{defn}[சமச்சீர்ப் பண்பு]
$Rxy$ ஆக இருக்கும்போதெல்லாம் $Ryx$ என்பதும் உண்மையாக இருக்கும்போது,
அப்போதும் அப்போது மட்டுமே $R \subseteq A^2$ என்னும் தொடர்பு
\emph{சமச்சீர்ப் பண்புடையது} எனப்படும்.
\end{defn}

% SEGMENT OLP-0014-S05
\begin{defn}[எதிர்சமச்சீர்ப் பண்பு]
$Rxy$, $Ryx$ ஆகிய இரண்டும் உண்மையாக இருக்கும்போதெல்லாம் $x=y$
என இருக்கும்போது, அப்போதும் அப்போது மட்டுமே $R \subseteq A^2$
என்னும் தொடர்பு \emph{எதிர்\-சமச்\-சீர்ப் பண்புடையது} எனப்படும்
(வேறுவிதமாகச் சொன்னால், $x\neq y$ எனில் $\lnot Rxy$ அல்லது $\lnot Ryx$).
\end{defn}

% SEGMENT OLP-0014-S06
\begin{explain}
சமச்சீர்த் தொடர்பில் $Rxy$, $Ryx$ ஆகியவை எப்போதும் ஒன்றாக உண்மையாகும்,
அல்லது இரண்டும் பொய்யாகும். எதிர்சமச்சீர்த் தொடர்பில்,
$Rxy$, $Ryx$ ஆகிய இரண்டும் ஒரே நேரத்தில் உண்மையாக இருக்க வேண்டுமெனில்,
$x = y$ ஆக இருக்க வேண்டும்.
ஆனால் $x = y$ ஆக இருக்கும்போது $Rxy$, $Ryx$ ஆகியவை உண்மையாக
\emph{இருக்க வேண்டும்} என இது கோருவதில்லை;
அவை உண்மையாக இருப்பதை இது விலக்கவில்லை என்பதே பொருள்.
எனவே எதிர்சமச்சீர்த் தொடர்பு தற்சுட்டுப் பண்புடையதாக இருக்கலாம்;
ஆனால் ஒவ்வோர் எதிர்சமச்சீர்த் தொடர்பும் தற்சுட்டுப் பண்புடையதல்ல.
எதிர்சமச்சீராக இருப்பதும், வெறுமனே சமச்சீராக இல்லாதிருப்பதும்
வெவ்வேறு நிபந்தனைகள் என்பதையும் கவனிக்கவும்.
உண்மையில், ஒரு தொடர்பு ஒரே நேரத்தில் சமச்சீராகவும் எதிர்சமச்சீராகவும்
இருக்கலாம் (எடுத்துக்காட்டாக, ஒன்றாதல் தொடர்பு அத்தகையது).
\end{explain}

% SEGMENT OLP-0014-S07
\begin{defn}[ஒப்பிடத்தக்க பண்பு]
அனைத்து $x,y\in A$ க்கும், $x \neq y$ எனில் $Rxy$ அல்லது $Ryx$
ஆக இருந்தால், $R \subseteq A^2$ என்னும் தொடர்பு
\emph{ஒப்பிடத்தக்க பண்புடையது} எனப்படும்.
\end{defn}

% SEGMENT OLP-0014-S08
\begin{prob}
பின்வரும் பண்புகளைக் கொண்ட தொடர்புகளுக்கு எடுத்துக்காட்டுகள் தருக:
(a) தற்சுட்டு, சமச்சீர் ஆகிய பண்புகள் உள்ளன, ஆனால் கடப்புப் பண்பு இல்லை;
(b) தற்சுட்டு, எதிர்சமச்சீர் ஆகிய பண்புகள் உள்ளன;
(c) எதிர்சமச்சீர், கடப்பு ஆகிய பண்புகள் உள்ளன, ஆனால் தற்சுட்டுப் பண்பு இல்லை;
(d) தற்சுட்டு, சமச்சீர், கடப்பு ஆகிய மூன்று பண்புகளும் உள்ளன.
எண்கள் அல்லது கணங்களின் மீதான தொடர்புகளைப் பயன்படுத்த வேண்டாம்.
\end{prob}

% SEGMENT OLP-0014-S09
\begin{defn}[எங்கும் தற்சுட்டின்மை]
அனைத்து $x \in A$ க்கும் $Rxx$ என்பது பொய்யாக இருந்தால்,
$R \subseteq A^2$ என்னும் தொடர்பு \emph{எங்கும் தற்சுட்டற்றது} எனப்படும்.
\end{defn}

% SEGMENT OLP-0014-S10
\begin{defn}[முழுச் சமச்சீரின்மை]
எந்த ஒரு ஜோடி $x,y\in A$ க்கும் $Rxy$, $Ryx$ ஆகிய இரண்டும்
உண்மையாக இருப்பதில்லை எனில், $R \subseteq A^2$ என்னும் தொடர்பு
\emph{முழுச் சமச்சீரின்மை உடையது} எனப்படும்.
\end{defn}

% SEGMENT OLP-0014-S11
$A \neq \emptyset$ எனில், $A$ இன் மீதான எங்கும் தற்சுட்டற்ற எந்தத்
தொடர்பும் தற்சுட்டுப் பண்புடையதாக இருக்காது என்பதைக் கவனிக்கவும்.
மேலும், $A$ இன் மீதான முழுச் சமச்சீரின்மை உடைய ஒவ்வொரு தொடர்பும்
எதிர்சமச்சீர்ப் பண்புடையதாகவும் இருக்கும்.
எனினும், தற்சுட்டுப் பண்பும் இல்லாமல் எங்கும் தற்சுட்டற்றதாகவும் இல்லாத
$R \subseteq A^2$ என்ற தொடர்புகள் உள்ளன.
முழுச் சமச்சீரின்மை இல்லாத எதிர்சமச்சீர்த் தொடர்புகளும் உள்ளன.

\end{document}
